\documentclass[11pt,a4paper]{article}

% ------------------------------------------------------------------
%  Modelos Avanzados de Computacion - Relacion 4 (resuelta)
%  Resolucion con fines de estudio.
% ------------------------------------------------------------------

\usepackage[utf8]{inputenc}
\usepackage[T1]{fontenc}
\usepackage{babel}
\babelprovide[import, main]{spanish}
\usepackage{amsmath,amssymb,amsthm}
\usepackage{mathtools}
\usepackage{enumitem}
\usepackage{geometry}
\usepackage{xcolor}
\usepackage{tikz}
\usetikzlibrary{arrows.meta,positioning,automata}
\usepackage{fancyhdr}
\usepackage{hyperref}
\hypersetup{colorlinks=true,linkcolor=blue,urlcolor=blue}

\geometry{margin=2.4cm}

% --- Entornos de teoremas -----------------------------------------
\theoremstyle{plain}
\newtheorem{teorema}{Teorema}
\newtheorem{lema}{Lema}
\newtheorem{proposicion}{Proposici\'on}
\theoremstyle{definition}
\newtheorem{definicion}{Definici\'on}
\newtheorem*{nota}{Nota}
\newtheorem*{idea}{Idea}

% --- Macros utiles ------------------------------------------------
\newcommand{\NP}{\mathrm{NP}}
\newcommand{\coNP}{\mathrm{CoNP}}
\newcommand{\Pclass}{\mathrm{P}}
\newcommand{\Lclass}{\mathrm{L}}
\newcommand{\SPACE}{\mathrm{ESPACIO}}
\newcommand{\TIME}{\mathrm{TIEMPO}}
\newcommand{\reduce}{\propto}
\newcommand{\Lstar}{L^{*}}
\renewcommand{\qedsymbol}{$\blacksquare$}

% --- Encabezados --------------------------------------------------
\pagestyle{fancy}
\fancyhf{}
\lhead{Modelos Avanzados de Computaci\'on}
\rhead{Relaci\'on 4}
\cfoot{\thepage}

\title{\textbf{Modelos Avanzados de Computaci\'on}\\[2mm]
\large Relaci\'on 4 --- Resoluci\'on para estudio}
\author{}
\date{}

\begin{document}
\maketitle
\thispagestyle{fancy}

\begin{nota}
Documento de estudio con las soluciones de la Relaci\'on 4. La relaci\'on cubre
clases de complejidad ($\Pclass$, $\NP$, $\coNP$, $\Lclass=\SPACE(\log n)$,
$\SPACE(n)$), reducciones en espacio logar\'itmico, grafos y composici\'on de
funciones. Las t\'ecnicas empleadas (verificadores con certificado, simulaci\'on
con punteros en espacio $O(\log n)$, argumentos de cierre por uni\'on/intersecci\'on,
diagonalizaci\'on impl\'icita v\'ia teoremas de jerarqu\'ia) son resultados
cl\'asicos del \'area, y coinciden con los desarrollos est\'andar de Sipser
(\emph{Introduction to the Theory of Computation}), Papadimitriou
(\emph{Computational Complexity}) y Garey--Johnson (\emph{Computers and
Intractability}). Donde un argumento sigue de cerca uno de estos textos se
indica expl\'icitamente.
\end{nota}

\tableofcontents
\vspace{4mm}
\hrule
\vspace{4mm}

% ==================================================================
\section{Grafos dirigidos ac\'iclicos (DAG)}
% ==================================================================

\subsection*{(a) Todo DAG finito tiene una fuente}

\begin{proposicion}
Todo grafo dirigido finito y ac\'iclico posee al menos una \emph{fuente}: un nodo
sin arcos entrantes.
\end{proposicion}

\begin{proof}
Razonamos por contradicci\'on. Sea $G=(V,E)$ finito y ac\'iclico, y supongamos que
\emph{ning\'un} nodo es fuente; es decir, todo nodo tiene al menos un arco entrante.

Construimos un camino hacia atr\'as. Tomamos un nodo cualquiera $v_0$. Como $v_0$ no
es fuente, existe un arco $(v_1,v_0)\in E$. Como $v_1$ tampoco es fuente, existe
$(v_2,v_1)\in E$. Repitiendo, generamos una sucesi\'on
\[
\cdots \to v_3 \to v_2 \to v_1 \to v_0,
\]
en la que cada $v_{i+1}$ tiene un arco hacia $v_i$. Como $|V|=n$ es finito, entre
los $n+1$ primeros nodos $v_0,v_1,\dots,v_n$ debe haber, por el principio del
palomar, dos \'indices $i<j$ con $v_i=v_j$. El tramo
\[
v_j \to v_{j-1}\to \cdots \to v_i = v_j
\]
es entonces un \textbf{ciclo dirigido}, lo que contradice que $G$ sea ac\'iclico.

Por tanto, alg\'un nodo no puede tener arco entrante: $G$ tiene una fuente.
\end{proof}

\begin{nota}
Por simetr\'ia (invirtiendo todos los arcos) el mismo argumento prueba que todo DAG
finito tiene tambi\'en un \emph{sumidero} (nodo sin arcos salientes). Esto se usar\'a
impl\'icitamente para el orden topol\'ogico.
\end{nota}

\subsection*{(b) Numeraci\'on topol\'ogica}

\begin{proposicion}
Un grafo dirigido con $n$ nodos es ac\'iclico si y solo si se pueden numerar sus nodos
$1,2,\dots,n$ de modo que todo arco vaya de un n\'umero menor a uno mayor.
\end{proposicion}

\begin{proof}
$(\Leftarrow)$ Supongamos que existe tal numeraci\'on $\nu:V\to\{1,\dots,n\}$ con
$\nu(u)<\nu(v)$ para todo arco $(u,v)$. Si hubiera un ciclo
$w_1\to w_2\to\cdots\to w_k\to w_1$, encadenando las desigualdades obtendr\'iamos
\[
\nu(w_1)<\nu(w_2)<\cdots<\nu(w_k)<\nu(w_1),
\]
es decir $\nu(w_1)<\nu(w_1)$, absurdo. Luego $G$ es ac\'iclico.

$(\Rightarrow)$ Supongamos $G$ ac\'iclico. Construimos la numeraci\'on por inducci\'on
sobre $n$ (este es el algoritmo de \emph{ordenaci\'on topol\'ogica}).

Por el apartado (a), $G$ tiene una fuente $s$. Le asignamos el n\'umero $1$:
$\nu(s)=1$. Consideramos $G'=G\setminus\{s\}$ (eliminamos $s$ y sus arcos
salientes). $G'$ sigue siendo ac\'iclico (quitar nodos/arcos no crea ciclos) y tiene
$n-1$ nodos; por hip\'otesis de inducci\'on admite una numeraci\'on
$\nu':V\setminus\{s\}\to\{1,\dots,n-1\}$ compatible. Definimos
$\nu(v)=\nu'(v)+1$ para $v\neq s$.

Comprobamos que todo arco va de menor a mayor:
\begin{itemize}[nosep]
  \item Arcos que salen de $s$: van de $\nu(s)=1$ a alg\'un $\nu(v)\geq 2$. Correcto.
  \item Arcos que no tocan $s$: estaban en $G'$ y eran crecientes para $\nu'$, luego
        siguen siendo crecientes para $\nu=\nu'+1$. Correcto.
  \item No hay arcos \emph{entrantes} a $s$ porque $s$ es fuente.
\end{itemize}
Esto completa la inducci\'on.
\end{proof}

\subsection*{(c) Algoritmo polin\'omico para decidir si un grafo es ac\'iclico}

La prueba de (b) es constructiva y da directamente el algoritmo (ordenaci\'on
topol\'ogica mediante eliminaci\'on iterativa de fuentes, conocido como
algoritmo de Kahn):

\begin{quote}
\textbf{Algoritmo} \textsc{EsAciclico}$(G=(V,E))$:
\begin{enumerate}[nosep]
  \item Calcular el grado de entrada $d^{-}(v)$ de cada nodo. Inicializar una cola
        $Q$ con todos los nodos de grado de entrada $0$ (las fuentes). Contador
        $c\leftarrow 0$.
  \item Mientras $Q\neq\emptyset$:
        \begin{enumerate}[nosep]
          \item Extraer $u$ de $Q$; $c\leftarrow c+1$.
          \item Para cada arco $(u,w)$: decrementar $d^{-}(w)$; si queda en $0$,
                a\~nadir $w$ a $Q$.
        \end{enumerate}
  \item Si $c=|V|$ devolver \textsc{Aciclico}; en otro caso devolver
        \textsc{Tiene Ciclo}.
\end{enumerate}
\end{quote}

\paragraph{Correcci\'on.} Si $G$ es ac\'iclico, por (a) en cada paso del subgrafo
restante siempre queda una fuente, as\'i que el proceso elimina los $n$ nodos y
$c=n$. Rec\'iprocamente, si el proceso se detiene con $c<n$, todos los nodos
restantes tienen grado de entrada $\geq 1$, luego (por el argumento de (a)
aplicado al subgrafo restante) ese subgrafo contiene un ciclo, y por tanto $G$
tambi\'en.

\paragraph{Complejidad.} Cada arco se procesa exactamente una vez (al eliminar su
origen) y cada nodo se encola/desencola una vez. El coste es
$O(|V|+|E|)$, lineal y en particular polin\'omico.

% ==================================================================
\section{Grafos bipartitos}
% ==================================================================

\subsection*{(a) Bipartito $\Leftrightarrow$ todos los ciclos de longitud par}

Trabajamos con grafos \emph{no dirigidos}.

\begin{proposicion}
Un grafo $G=(V,E)$ es bipartito si y solo si todos sus ciclos tienen longitud par.
\end{proposicion}

\begin{proof}
$(\Rightarrow)$ Supongamos $G$ bipartito con partici\'on $V=A\,\dot\cup\,B$ (toda
arista une un v\'ertice de $A$ con uno de $B$). Sea
$v_0,v_1,\dots,v_k=v_0$ un ciclo. Cada arista cambia de lado: si $v_0\in A$ entonces
$v_1\in B$, $v_2\in A$, \dots; en general $v_i\in A$ si $i$ es par y $v_i\in B$ si $i$
es impar. Como debe cerrarse en $v_k=v_0\in A$, el \'indice $k$ ha de ser par. Luego
todo ciclo tiene longitud par.

$(\Leftarrow)$ Supongamos que todos los ciclos son de longitud par. Basta trabajar
por componentes conexas (la uni\'on de biparticiones de cada componente es una
biparticion del grafo). Fijada una componente, elegimos una ra\'iz $r$ y
definimos, para cada v\'ertice $v$, la distancia $d(v)=$ longitud del camino m\'as
corto de $r$ a $v$. Coloreamos:
\[
A=\{v: d(v)\text{ par}\},\qquad B=\{v: d(v)\text{ impar}\}.
\]
Veamos que ninguna arista une dos v\'ertices del mismo color. Sea $\{u,v\}\in E$.
Si fuera $d(u)\equiv d(v)\pmod 2$, considere\-mos los caminos m\'as cortos
$P_u$ de $r$ a $u$ y $P_v$ de $r$ a $v$. Sea $w$ su \'ultimo v\'ertice com\'un. Los
tramos $w\!\to\!u$ y $w\!\to\!v$ tienen longitudes $d(u)-d(w)$ y $d(v)-d(w)$, de la
\emph{misma} paridad. El recorrido
\[
w\to\cdots\to u \;-\; v\to\cdots\to w
\]
(bajar a $u$, usar la arista $\{u,v\}$, subir desde $v$) es un ciclo de longitud
\[
\big(d(u)-d(w)\big) + 1 + \big(d(v)-d(w)\big),
\]
que es \textbf{impar} (suma de dos n\'umeros de igual paridad m\'as $1$). Esto
contradice la hip\'otesis. Por tanto toda arista une colores distintos y la
componente es bipartita.
\end{proof}

\subsection*{(b) Algoritmo polin\'omico para comprobar bipartici\'on}

La prueba sugiere directamente un algoritmo de \emph{2-coloreado} por BFS:

\begin{quote}
\textbf{Algoritmo} \textsc{EsBipartito}$(G)$:
\begin{enumerate}[nosep]
  \item Marcar todos los v\'ertices como sin color.
  \item Para cada v\'ertice $r$ a\'un sin color (cabecera de una componente):
        \begin{enumerate}[nosep]
          \item Asignar $color(r)=0$ y lanzar un BFS desde $r$.
          \item Al visitar una arista $\{u,v\}$: si $v$ no tiene color, asignar
                $color(v)=1-color(u)$; si ya lo tiene y $color(v)=color(u)$,
                devolver \textsc{No bipartito}.
        \end{enumerate}
  \item Si nunca se produjo conflicto, devolver \textsc{Bipartito} con la
        partici\'on dada por los colores.
\end{enumerate}
\end{quote}

\paragraph{Correcci\'on.} El BFS asigna a cada v\'ertice la paridad de su distancia
a la ra\'iz, exactamente el coloreado de la prueba de (a). Un conflicto
($color(u)=color(v)$ en una arista) corresponde a la situaci\'on del ciclo impar; su
ausencia certifica una biparticion v\'alida.

\paragraph{Complejidad.} Es un recorrido BFS est\'andar: $O(|V|+|E|)$, polin\'omico.

% ==================================================================
\section{$\Pclass$ es cerrada para uni\'on e intersecci\'on}
% ==================================================================

\begin{teorema}
Si $L_1,L_2\in\Pclass$, entonces $L_1\cup L_2\in\Pclass$ y
$L_1\cap L_2\in\Pclass$.
\end{teorema}

\begin{proof}
Por hip\'otesis existen m\'aquinas de Turing deterministas $M_1$ y $M_2$ que deciden
$L_1$ y $L_2$ en tiempo $O(n^{c_1})$ y $O(n^{c_2})$ respectivamente, para ciertas
constantes $c_1,c_2$.

Construimos una m\'aquina $M$ que, con entrada $x$ de longitud $n$:
\begin{enumerate}[nosep]
  \item Simula $M_1(x)$ y guarda su respuesta $b_1\in\{\text{S\'i},\text{No}\}$.
  \item Simula $M_2(x)$ y guarda su respuesta $b_2$.
  \item \textbf{Uni\'on:} acepta si $b_1=\text{S\'i}$ \emph{o} $b_2=\text{S\'i}$.\\
        \textbf{Intersecci\'on:} acepta si $b_1=\text{S\'i}$ \emph{y}
        $b_2=\text{S\'i}$.
\end{enumerate}

$M$ es determinista y siempre se detiene. Su tiempo es a lo sumo
\[
O(n^{c_1}) + O(n^{c_2}) + O(1) = O\!\big(n^{\max(c_1,c_2)}\big),
\]
es decir, polin\'omico. Adem\'as $L(M)=L_1\cup L_2$ (resp. $L_1\cap L_2$) por
construcci\'on. Luego ambos lenguajes est\'an en $\Pclass$.
\end{proof}

\begin{nota}
El mismo argumento prueba que $\Pclass$ es cerrada para complementaci\'on: basta
intercambiar los estados de aceptaci\'on y rechazo de una m\'aquina polin\'omica
determinista (la determinaci\'on y la parada garantizada son esenciales aqu\'i).
\end{nota}

% ==================================================================
\section{Cierre de clases de funciones por composici\'on polin\'omica}
% ==================================================================

\begin{definicion}[recordatorio del enunciado]
Sea $\mathcal{C}$ una clase de funciones $\mathbb{N}\to\mathbb{N}$ y $p$ un polinomio
arbitrario.
\begin{itemize}[nosep]
  \item $\mathcal{C}$ es cerrada para la composici\'on polin\'omica \textbf{por la
        izquierda} si $f\in\mathcal{C}\Rightarrow p(f(n))=O(g(n))$ para alguna
        $g\in\mathcal{C}$.
  \item $\mathcal{C}$ es cerrada para la composici\'on polin\'omica \textbf{por la
        derecha} si $f\in\mathcal{C}\Rightarrow f(p(n))=O(g(n))$ para alguna
        $g\in\mathcal{C}$.
\end{itemize}
Como indica la nota del enunciado, por tratarse de \'ordenes de complejidad basta
considerar $p(n)=n^{l}$ con $l\in\mathbb{N}$.
\end{definicion}

\begin{idea}
\emph{Por la izquierda} aplicamos primero $f$ y luego elevamos a una potencia:
$p(f(n))=f(n)^{l}$. \emph{Por la derecha} sustituimos $n$ por $n^{l}$ dentro de $f$:
$f(n^{l})$. La clase es cerrada si el resultado vuelve a estar dominado por una
funci\'on de la misma familia.
\end{idea}

A continuaci\'on analizamos las seis familias. Resumimos al final en una tabla.

% ---------------------------------------------------------------
\subsection*{(a) $\mathcal{C}=\{n^{k}: k>0\}$ (polinomios mon\'omicos)}

Sea $f(n)=n^{k}\in\mathcal{C}$.

\paragraph{Izquierda.} $p(f(n)) = (n^{k})^{l} = n^{kl}$. Como $kl>0$, la funci\'on
$n^{kl}\in\mathcal{C}$, y $p(f(n))=O(n^{kl})$. \textbf{Cerrada por la izquierda.}

\paragraph{Derecha.} $f(p(n)) = (n^{l})^{k} = n^{lk}\in\mathcal{C}$.
\textbf{Cerrada por la derecha.}

\medskip\noindent\fbox{(a) cerrada por ambos lados.}

% ---------------------------------------------------------------
\subsection*{(b) $\mathcal{C}=\{n\cdot k: k>0\}$ (funciones lineales)}

Aqu\'i $f(n)=k\,n$ es lineal con pendiente $k>0$ (el factor constante, no el
exponente).

\paragraph{Izquierda.} $p(f(n)) = (k\,n)^{l} = k^{l}\,n^{l}$. Para $l\geq 2$ esto es
$\Theta(n^{l})$, un crecimiento \emph{superlineal} que \textbf{no} puede acotarse por
ninguna funci\'on lineal $g(n)=k'n$ (pues $n^{l}/n=n^{l-1}\to\infty$). \textbf{No
cerrada por la izquierda.}

\paragraph{Derecha.} $f(p(n)) = k\,(n^{l}) = k\,n^{l}$. Igual que antes, para
$l\geq 2$ es $\Theta(n^{l})$, no acotable por una funci\'on lineal. \textbf{No
cerrada por la derecha.}

\medskip\noindent\fbox{(b) no cerrada por ning\'un lado.}

% ---------------------------------------------------------------
\subsection*{(c) $\mathcal{C}=\{k^{n}: k>0\}$ (exponenciales de base constante)}

Sea $f(n)=k^{n}$. (Tomamos $k\geq 2$; para $k=1$ es la constante $1$ y para $k>1$ el
an\'alisis es el interesante.)

\paragraph{Izquierda.} $p(f(n)) = (k^{n})^{l} = k^{ln} = (k^{l})^{n}$. Llamando
$k'=k^{l}>0$ obtenemos $(k')^{n}\in\mathcal{C}$. \textbf{Cerrada por la izquierda.}

\paragraph{Derecha.} $f(p(n)) = k^{\,n^{l}}$. Para $l\geq 2$ esto crece como
$k^{n^{2}}$, mucho m\'as r\'apido que cualquier $(k')^{n}=k'^{\,n}$, ya que
\[
\frac{k^{\,n^{l}}}{(k')^{\,n}} = \frac{k^{\,n^{l}}}{k^{\,n\log_k k'}}
= k^{\,n^{l}-n\log_k k'}\longrightarrow\infty.
\]
No existe $g\in\mathcal{C}$ que lo domine. \textbf{No cerrada por la derecha.}

\medskip\noindent\fbox{(c) cerrada por la izquierda, no por la derecha.}

% ---------------------------------------------------------------
\subsection*{(d) $\mathcal{C}=\{2^{\,n^{k}}: k>0\}$ (exponenciales con exponente polin\'omico)}

Sea $f(n)=2^{\,n^{k}}$.

\paragraph{Izquierda.} $p(f(n)) = \big(2^{\,n^{k}}\big)^{l} = 2^{\,l\,n^{k}}$.
Necesitamos acotar $2^{\,l\,n^{k}}$ por alg\'un $2^{\,n^{k'}}$. Como $l$ es
constante, para $n$ grande $l\,n^{k}\leq n^{k+1}$, luego
\[
2^{\,l\,n^{k}} \leq 2^{\,n^{k+1}} = O\!\big(2^{\,n^{k+1}}\big),\qquad
2^{\,n^{k+1}}\in\mathcal{C}.
\]
\textbf{Cerrada por la izquierda.}

\paragraph{Derecha.} $f(p(n)) = 2^{\,(n^{l})^{k}} = 2^{\,n^{lk}}$. Como $lk>0$, la
funci\'on $2^{\,n^{lk}}\in\mathcal{C}$. \textbf{Cerrada por la derecha.}

\medskip\noindent\fbox{(d) cerrada por ambos lados.}

% ---------------------------------------------------------------
\subsection*{(e) $\mathcal{C}=\{\log^{k}(n): k>0\}$ (potencias de logaritmo)}

Sea $f(n)=(\log n)^{k}=\log^{k} n$ (con redondeo al entero m\'as cercano, irrelevante
para el orden).

\paragraph{Izquierda.} $p(f(n)) = (\log^{k} n)^{l} = \log^{kl} n\in\mathcal{C}$, pues
$kl>0$. \textbf{Cerrada por la izquierda.}

\paragraph{Derecha.} $f(p(n)) = \log^{k}(n^{l}) = (l\log n)^{k} = l^{k}\,\log^{k} n
= O(\log^{k} n)$, y $\log^{k} n\in\mathcal{C}$. \textbf{Cerrada por la derecha.}

\medskip\noindent\fbox{(e) cerrada por ambos lados.}

% ---------------------------------------------------------------
\subsection*{(f) $\mathcal{C}=\{\log n\}$ (un \'unico elemento: el logaritmo)}

Ahora la clase tiene \emph{una sola} funci\'on, $f(n)=\log n$, as\'i que la \'unica
candidata posible para $g$ es la propia $\log n$.

\paragraph{Izquierda.} $p(f(n)) = (\log n)^{l} = \log^{l} n$. Para $l\geq 2$ se tiene
$\log^{l} n / \log n = \log^{l-1} n\to\infty$, de modo que $\log^{l} n\neq O(\log n)$.
La \'unica $g$ disponible no lo acota. \textbf{No cerrada por la izquierda.}

\paragraph{Derecha.} $f(p(n)) = \log(n^{l}) = l\log n = O(\log n)$, y la \'unica
funci\'on de la clase, $\log n$, s\'i lo domina. \textbf{Cerrada por la derecha.}

\medskip\noindent\fbox{(f) cerrada por la derecha, no por la izquierda.}

% ---------------------------------------------------------------
\subsection*{Resumen}

\begin{center}
\renewcommand{\arraystretch}{1.35}
\begin{tabular}{|l|c|c|}
\hline
\textbf{Clase} & \textbf{Izquierda} $p(f(n))=f(n)^{l}$ & \textbf{Derecha} $f(n^{l})$\\
\hline
(a) $\{n^{k}\}$        & S\'i & S\'i \\
(b) $\{k\,n\}$         & No   & No   \\
(c) $\{k^{n}\}$        & S\'i & No   \\
(d) $\{2^{\,n^{k}}\}$  & S\'i & S\'i \\
(e) $\{\log^{k} n\}$   & S\'i & S\'i \\
(f) $\{\log n\}$       & No   & S\'i \\
\hline
\end{tabular}
\end{center}

\begin{nota}
Patr\'on conceptual: la composici\'on \emph{por la izquierda} eleva $f$ a la
potencia $l$; sobrevive cuando la clase es estable al elevar a potencias
constantes (familias param\'etricas en el exponente o en el \'indice de
logaritmo). La composici\'on \emph{por la derecha} mete $n^{l}$ dentro de $f$;
sobrevive cuando ``elevar el argumento a una potencia'' equivale a moverse dentro de
la familia. Las familias con \emph{un solo grado de libertad mal situado} (lineales,
o el logaritmo aislado) fallan en el lado que ataca ese grado de libertad.
\end{nota}

% ==================================================================
\section{$L=\{wcw : w\in\{0,1\}^{*}\}$ se reconoce en espacio $\log n$}
% ==================================================================

\begin{teorema}
El lenguaje $L=\{wcw : w\in\{0,1\}^{*}\}$ pertenece a
$\Lclass=\SPACE(\log n)$.
\end{teorema}

\begin{idea}
No podemos copiar $w$ a la cinta de trabajo (eso costar\'ia espacio lineal). En su
lugar usamos \emph{punteros}: comparamos las dos mitades posici\'on a posici\'on
guardando solo \'indices, y cada \'indice cabe en $O(\log n)$ bits.
\end{idea}

\begin{proof}
Trabajamos con el modelo est\'andar de m\'aquina de Turing con cinta de entrada de
\emph{solo lectura} y una cinta de trabajo aparte cuyo espacio es el que contamos.
Sea la entrada de longitud $n$.

\textbf{Fase 1: validar el formato.} Recorremos la entrada y comprobamos que contiene
\emph{exactamente} un s\'imbolo $c$. Para ello mantenemos un contador del n\'umero de
$c$ vistos y la posici\'on $m$ de la $c$. El contador y la posici\'on son n\'umeros
$\leq n$, luego ocupan $O(\log n)$ bits. Si el n\'umero de $c$ no es $1$, rechazar.
Sea $m$ la posici\'on de la $c$ (separa la entrada en parte izquierda
$u=x_1\cdots x_{m-1}$ y parte derecha $v=x_{m+1}\cdots x_{n}$).

\textbf{Fase 2: comprobar longitudes iguales.} La pertenencia exige $|u|=|v|$, es
decir $m-1 = n-m$. Esto se verifica con aritm\'etica sobre los \'indices $m$ y $n$,
de nuevo en $O(\log n)$ espacio. Si no coinciden, rechazar.

\textbf{Fase 3: comparar car\'acter a car\'acter.} Para
$i=1,2,\dots,m-1$ comparamos el s\'imbolo en la posici\'on $i$ (mitad izquierda) con
el de la posici\'on $m+i$ (mitad derecha). Concretamente, mantenemos un contador $i$
y, en cada iteraci\'on:
\begin{itemize}[nosep]
  \item movemos la cabeza de \emph{solo lectura} a la posici\'on $i$ y leemos
        $x_i$ (guardando ese \'unico bit);
  \item movemos la cabeza a la posici\'on $m+i$ y leemos $x_{m+i}$;
  \item si $x_i\neq x_{m+i}$, rechazar.
\end{itemize}
Si todas las comparaciones coinciden, aceptar.

\textbf{An\'alisis de espacio.} En la cinta de trabajo solo almacenamos en cada
momento un contador (el \'indice $i$, o $m$, o el n\'umero de $c$) y a lo sumo uno o
dos bits le\'idos. Todos los \'indices est\'an acotados por $n$, luego se codifican en
$\lceil\log_2 n\rceil + O(1)$ bits. El n\'umero de variables es constante. Por tanto
el espacio total es $O(\log n)$.

(El tiempo es polin\'omico --- $O(n^2)$ por los reposicionamientos de la cabeza ---
pero el enunciado solo exige la cota de espacio.)
\end{proof}

% ==================================================================
\section{Si $L\in\Pclass$ entonces $\Lstar\in\Pclass$}
% ==================================================================

\begin{teorema}
$\Pclass$ es cerrada para la estrella de Kleene: si $L\in\Pclass$, entonces
$\Lstar\in\Pclass$.
\end{teorema}

\begin{idea}
$x\in\Lstar$ significa que $x$ se puede \emph{partir} en bloques consecutivos
$x=w_1 w_2\cdots w_k$ con cada $w_j\in L$ (y $\varepsilon\in\Lstar$ siempre). Esto se
decide con programaci\'on din\'amica sobre los prefijos de $x$.
\end{idea}

\begin{proof}
Sea $M$ una m\'aquina que decide $L$ en tiempo $O(n^{c})$. Dada
$x=x_1 x_2\cdots x_n$, escribimos $x[i..j]=x_i\cdots x_j$ (subcadena), y por convenio
$x[i..i-1]=\varepsilon$.

Definimos un vector booleano $A[0..n]$ con la interpretaci\'on
\[
A[j] = \text{Verdadero} \iff \text{el prefijo } x[1..j]\in\Lstar .
\]
Recurrencia (descomponiendo seg\'un d\'onde empieza el \'ultimo bloque):
\[
A[0]=\text{Verdadero (el prefijo vac\'io est\'a en }\Lstar),
\]
\[
A[j]=\bigvee_{0\le i<j}\Big(A[i]\ \wedge\ \big(x[i+1..j]\in L\big)\Big),
\quad 1\le j\le n .
\]
Es decir, $x[1..j]\in\Lstar$ si existe un punto de corte $i$ tal que el prefijo
$x[1..i]$ ya estaba en $\Lstar$ y el \'ultimo trozo $x[i+1..j]$ pertenece a $L$.

\textbf{Algoritmo.} Calcular $A[0],A[1],\dots,A[n]$ en orden creciente usando la
recurrencia; para cada par $(i,j)$ se invoca $M$ sobre la subcadena $x[i+1..j]$.
Aceptar si $A[n]=\text{Verdadero}$.

\textbf{Correcci\'on.} Por inducci\'on sobre $j$: $A[j]$ es verdadero exactamente
cuando $x[1..j]$ admite una descomposici\'on en factores de $L$, que es justo la
definici\'on de pertenencia a $\Lstar$. El caso base $A[0]$ corresponde a
$\varepsilon\in\Lstar$. En particular $A[n]$ decide $x\in\Lstar$.

\textbf{Complejidad.} Hay $O(n^{2})$ pares $(i,j)$. Cada uno requiere una llamada a
$M$ sobre una subcadena de longitud $\le n$, con coste $O(n^{c})$. El total es
\[
O(n^{2})\cdot O(n^{c}) = O(n^{c+2}),
\]
polin\'omico. Por tanto $\Lstar\in\Pclass$.
\end{proof}

% ==================================================================
\section{Si $L\in\NP$ entonces $\Lstar\in\NP$}
% ==================================================================

\begin{teorema}
$\NP$ es cerrada para la estrella de Kleene: si $L\in\NP$, entonces $\Lstar\in\NP$.
\end{teorema}

\begin{idea}
Caracterizamos $\NP$ por verificador y certificado. Para $\Lstar$ el certificado
consta de: (i) los \emph{puntos de corte} que parten $x$ en bloques, y (ii) un
certificado de pertenencia a $L$ para cada bloque. Todo ello es de tama\~no
polin\'omico y verificable en tiempo polin\'omico.
\end{idea}

\begin{proof}[Demostraci\'on (versi\'on verificador)]
Como $L\in\NP$, existe un verificador determinista $V$ y un polinomio $q$ tales que,
para toda cadena $w$,
\[
w\in L \iff \exists\, c_w,\ |c_w|\le q(|w|),\ \ V(w,c_w)=1,
\]
y $V$ corre en tiempo polin\'omico.

Definimos un verificador $V^{*}$ para $\Lstar$. Dado $x$ de longitud $n$, el
certificado propuesto es
\[
\big(\, i_0=0 < i_1 < i_2 < \cdots < i_k = n,\ \ c_1, c_2,\dots, c_k \,\big),
\]
donde los $i_j$ marcan los cortes (el bloque $j$-\'esimo es $w_j=x[i_{j-1}+1..i_j]$)
y $c_j$ es un certificado de que $w_j\in L$. El caso $x=\varepsilon$ se acepta
directamente ($k=0$).

$V^{*}$ procede as\'i:
\begin{enumerate}[nosep]
  \item Comprobar que $0=i_0<i_1<\cdots<i_k=n$ (los cortes son v\'alidos y cubren
        todo $x$ sin solapes).
  \item Para cada $j=1,\dots,k$: extraer $w_j=x[i_{j-1}+1..i_j]$ y verificar
        $V(w_j,c_j)=1$.
  \item Aceptar si y solo si todas las verificaciones tienen \'exito.
\end{enumerate}

\textbf{Correcci\'on.} $x\in\Lstar$ syss existe una partici\'on de $x$ en bloques
$w_1\cdots w_k$ con cada $w_j\in L$, lo cual equivale a la existencia de cortes
v\'alidos y de certificados $c_j$ aceptados por $V$. Esto es exactamente lo que
$V^{*}$ comprueba.

\textbf{Tama\~no del certificado.} Hay a lo sumo $k\le n$ cortes, cada uno un \'indice
$\le n$: $O(n\log n)$ bits. Cada $c_j$ tiene tama\~no $\le q(|w_j|)\le q(n)$, y
$\sum_j |w_j|=n$, luego la suma de longitudes de los certificados es a lo sumo
$n\cdot q(n)$, polin\'omica.

\textbf{Tiempo de verificaci\'on.} Comprobar los cortes es lineal; ejecutar $V$
sobre cada bloque cuesta polin\'omico y se hace $k\le n$ veces. Total polin\'omico.

Por tanto $\Lstar\in\NP$.
\end{proof}

\begin{nota}[demostraci\'on alternativa con MT no determinista]
De forma equivalente: una m\'aquina no determinista para $\Lstar$ va leyendo $x$ de
izquierda a derecha, ``adivina'' d\'onde termina cada bloque, y para cada bloque
ejecuta (no deterministamente) la m\'aquina de $L$. Como $\sum|w_j|=n$, el c\'omputo
total es de longitud polin\'omica. Tambi\'en se puede usar la t\'ecnica de
programaci\'on din\'amica del Ejercicio 6 sustituyendo ``$x[i+1..j]\in L$'' por una
llamada al verificador no determinista; el esquema es id\'entico.
\end{nota}

% ==================================================================
\section{$\NP$ es cerrada para uni\'on e intersecci\'on}
% ==================================================================

\begin{teorema}
Si $L_1,L_2\in\NP$, entonces $L_1\cup L_2\in\NP$ y $L_1\cap L_2\in\NP$.
\end{teorema}

\begin{proof}
Usamos la caracterizaci\'on por verificadores. Sean $V_1,V_2$ verificadores
polin\'omicos y $q_1,q_2$ polinomios con
\[
x\in L_i \iff \exists\, c_i,\ |c_i|\le q_i(|x|),\ V_i(x,c_i)=1
\qquad (i=1,2).
\]

\paragraph{Uni\'on.} El certificado es un par (etiqueta, testigo): un bit
$b\in\{1,2\}$ que indica ``a qu\'e lenguaje pertenece $x$'' junto con el certificado
correspondiente $c_b$. El verificador $V_\cup(x,(b,c))$ ejecuta $V_b(x,c)$ y acepta
si este acepta. Entonces
\[
x\in L_1\cup L_2 \iff (\exists c_1: V_1(x,c_1)=1)\ \text{o}\ (\exists c_2:
V_2(x,c_2)=1) \iff \exists (b,c): V_\cup(x,(b,c))=1 .
\]
El certificado tiene tama\~no $1+\max(q_1,q_2)(|x|)$ y $V_\cup$ es polin\'omico. Luego
$L_1\cup L_2\in\NP$.

\paragraph{Intersecci\'on.} El certificado es el \emph{par} $(c_1,c_2)$ de ambos
testigos. El verificador $V_\cap(x,(c_1,c_2))$ ejecuta $V_1(x,c_1)$ y $V_2(x,c_2)$ y
acepta si \emph{ambos} aceptan. Entonces
\[
x\in L_1\cap L_2 \iff (\exists c_1: V_1(x,c_1)=1)\ \text{y}\ (\exists c_2:
V_2(x,c_2)=1)\iff \exists(c_1,c_2): V_\cap(x,(c_1,c_2))=1 .
\]

\medskip
Un detalle sobre la intersecci\'on: el cuantificador conjunto
``$\exists(c_1,c_2)$'' es correcto porque los dos testigos son
\emph{independientes} --- elegir uno no condiciona al otro --- de modo que
``$\exists c_1\,\exists c_2$'' equivale a ``existe el par''. El certificado tiene
tama\~no $q_1(|x|)+q_2(|x|)$, polin\'omico, y $V_\cap$ es polin\'omico. Luego
$L_1\cap L_2\in\NP$.
\end{proof}

% ==================================================================
\section{Si $\NP\neq\coNP$ entonces $\Pclass\neq\NP$}
% ==================================================================

\begin{teorema}
Si $\NP\neq\coNP$, entonces $\Pclass\neq\NP$.
\end{teorema}

\begin{idea}
Demostramos el contrarrec\'iproco: si $\Pclass=\NP$ entonces $\NP=\coNP$. La clave es
que $\Pclass$ es cerrada para complementaci\'on (Ejercicio 3 / Nota), propiedad que
\emph{no} se sabe que tenga $\NP$.
\end{idea}

\begin{proof}
Probamos la afirmaci\'on equivalente (contrarrec\'iproca):
\[
\Pclass=\NP \;\Longrightarrow\; \NP=\coNP .
\]

Supongamos $\Pclass=\NP$. Recordemos que $\Pclass$ es cerrada para complementaci\'on:
si $A\in\Pclass$, intercambiando aceptaci\'on y rechazo en una m\'aquina
determinista polin\'omica que decide $A$ obtenemos otra que decide $\overline{A}$ en
tiempo polin\'omico, luego $\overline{A}\in\Pclass$.

Veamos la doble inclusi\'on $\NP=\coNP$.

\textbf{$\NP\subseteq\coNP$:} Sea $A\in\NP$. Por hip\'otesis $A\in\Pclass$. Como
$\Pclass$ es cerrada para complemento, $\overline{A}\in\Pclass=\NP$. Pero
``$\overline{A}\in\NP$'' es precisamente la definici\'on de ``$A\in\coNP$''. Luego
$A\in\coNP$.

\textbf{$\coNP\subseteq\NP$:} Sea $A\in\coNP$, esto es $\overline{A}\in\NP$. Por
hip\'otesis $\overline{A}\in\Pclass$, y por cierre bajo complemento
$A\in\Pclass=\NP$. Luego $A\in\NP$.

De ambas inclusiones, $\NP=\coNP$.

Hemos probado $\Pclass=\NP\Rightarrow\NP=\coNP$. Tomando contrarrec\'iproco:
\[
\NP\neq\coNP \;\Longrightarrow\; \Pclass\neq\NP. \qedhere
\]
\end{proof}

% ==================================================================
\section{Par\'entesis bien emparejados est\'a en $\Lclass$}
% ==================================================================

\begin{teorema}
El lenguaje de las cadenas de par\'entesis $\{(,)\}$ correctamente emparejados y
anidados pertenece a $\Lclass=\SPACE(\log n)$.
\end{teorema}

\begin{idea}
Con un \emph{\'unico} tipo de par\'entesis no hace falta una pila: basta un
\emph{contador} de profundidad. La condici\'on de buen emparejamiento es local:
``el contador nunca se vuelve negativo y termina en cero''.
\end{idea}

\begin{proof}
Sea la entrada $x=x_1\cdots x_n$ sobre el alfabeto $\{(,)\}$. Mantenemos en la cinta
de trabajo un \'unico contador entero $b$ (``balance''), inicializado a $0$.
Recorremos la entrada de izquierda a derecha:
\begin{itemize}[nosep]
  \item si $x_i = $ \verb|(| , hacer $b\leftarrow b+1$;
  \item si $x_i = $ \verb|)| , hacer $b\leftarrow b-1$;
  \item si en alg\'un momento $b<0$, \textbf{rechazar} (hay un \verb|)| sin
        su \verb|(| previo).
\end{itemize}
Al terminar, \textbf{aceptar} si y solo si $b=0$.

\textbf{Correcci\'on.} Una cadena de un solo tipo de par\'entesis est\'a bien formada
syss en todo prefijo el n\'umero de \verb|(| es $\ge$ el de \verb|)| (nunca se cierra
de m\'as) y al final ambos coinciden (todo lo abierto se cierra). El contador $b$
mide exactamente \#\verb|(| $-$ \#\verb|)| del prefijo le\'ido; ``$b\ge 0$ siempre''
y ``$b=0$ al final'' son justo esas dos condiciones. As\'i,
\verb|(()())| y \verb|((()))| son aceptadas, mientras \verb|)()(| (el balance se
hace $-1$ al primer s\'imbolo) y \verb|(()))(| (balance llega a $-1$) son
rechazadas.

\textbf{Espacio.} El contador $b$ satisface $0\le b\le n$ en todo instante (mientras
no se rechace), luego se almacena en $\lceil\log_2(n+1)\rceil = O(\log n)$ bits. No se
usa nada m\'as. Por tanto el algoritmo es $O(\log n)$ en espacio, y el problema est\'a
en $\Lclass$.
\end{proof}

% ==================================================================
\section{Dos tipos de par\'entesis}
% ==================================================================

\begin{teorema}
El lenguaje de cadenas bien formadas con dos tipos de par\'entesis,
$(,)$ y $[,]$, tambi\'en pertenece a $\Lclass$.
\end{teorema}

\begin{idea}
Con dos tipos importa el \emph{orden} de cierre (anidamiento correcto:
\verb|([]())| s\'i, \verb|([)]| no). La forma natural de resolverlo es con una pila,
pero una pila puede crecer $\Theta(n)$. El truco para mantenernos en $O(\log n)$ es
\emph{no almacenar la pila}: para verificar cada cierre, recalculamos cu\'al es su
par\'entesis de apertura asociado recorriendo la entrada con contadores.
\end{idea}

\begin{proof}
\textbf{Condici\'on de validez.} La cadena es correcta syss:
\begin{enumerate}[nosep]
  \item est\'a balanceada \emph{ignorando el tipo} (el contador global de
        ``abiertos menos cerrados'' nunca es negativo y acaba en $0$); y
  \item cada s\'imbolo de cierre case en \emph{tipo} con su apertura asociada
        (la apertura que lo empareja en el anidamiento).
\end{enumerate}

\textbf{Emparejamiento sin pila.} Para una posici\'on de cierre $j$, su apertura
asociada es la posici\'on $i<j$ tal que el segmento $x[i+1..j-1]$ est\'a balanceado;
equivalentemente, $i$ es la \'ultima posici\'on $<j$ en la que el balance (contando
\emph{cualquier} tipo) tom\'o el valor $\big(\text{balance en }j\big)$ justo
\emph{antes} de cerrar. La idea operativa, en espacio $O(\log n)$:

Recorremos $x$ con un contador de balance $b$ (como en el Ejercicio 10, comprobando
de paso la condici\'on 1). Cada vez que encontramos un cierre en posici\'on $j$
(supongamos que deja el balance en valor $d=b$ tras decrementar), buscamos su
apertura: \emph{reiniciamos} un segundo recorrido desde el principio hasta $j$
llevando un contador auxiliar, y localizamos la \'ultima posici\'on $i<j$ que es una
apertura y desde la cual el sub-balance de $x[i+1..j-1]$ es $0$. Comprobamos que
$x_i$ y $x_j$ son del mismo tipo (ambos redondos o ambos cuadrados); si no, rechazar.

Todo el almacenamiento son unos pocos contadores acotados por $n$ (la posici\'on $j$,
la posici\'on candidata $i$, los balances), cada uno $O(\log n)$ bits. No guardamos la
pila completa: la ``recomputamos'' tantas veces como cierres haya. Esto sacrifica
tiempo (polin\'omico, del orden de $O(n^{2})$) a cambio de espacio.

\textbf{Espacio total:} $O(\log n)$. Por tanto el problema con dos tipos de
par\'entesis tambi\'en est\'a en $\Lclass$.

\medskip
\textbf{Segunda parte: lectura de izquierda a derecha sin retroceder requiere
$O(n)$.}

Si imponemos que la entrada se lee \emph{una sola vez, de izquierda a derecha, sin
volver atr\'as} (modelo \emph{online} / aut\'omata con memoria pero sin re-lectura),
entonces necesitamos recordar la secuencia de aperturas a\'un no cerradas, en orden:
es esencialmente una \textbf{pila} de tipos. Justificaci\'on (argumento de
distinguibilidad / \emph{fooling set}):

Consideremos los $2^{m}$ prefijos formados por $m$ aperturas cada una de tipo redondo
o cuadrado, p.\,ej. $u_S = s_1 s_2\cdots s_m$ con $s_t\in\{\,(\,,\,[\,\}$. Para dos
prefijos distintos $u_S\neq u_{S'}$ existe una primera posici\'on $t$ donde difieren
los tipos. La continuaci\'on que cierra en el orden correcto para $u_S$ ---es decir,
el ``espejo'' de tipos $\bar s_m\cdots\bar s_1$ donde $\bar{(}=)$ y $\bar{[}=]$---
completa $u_S$ a una cadena \emph{v\'alida}, pero aplicada a $u_{S'}$ produce un
\emph{desajuste de tipo} en alguna posici\'on, dando cadena \emph{inv\'alida}. Por
tanto la m\'aquina debe quedar en \emph{estados de memoria distintos} tras leer
$u_S$ y $u_{S'}$ (si coincidieran, aceptar\'ia o rechazar\'ia igual ambas
continuaciones, contradicci\'on).

As\'i, hay al menos $2^{m}$ configuraciones de memoria distinguibles tras leer $m$
s\'imbolos. Representar $2^{m}$ estados distintos exige al menos
$\log_2(2^{m}) = m$ bits de memoria. Como $m$ puede ser $\Theta(n)$, la memoria
necesaria es $\Omega(n)$, es decir $O(n)$ (lineal) en el peor caso.
\end{proof}

\begin{nota}
La diferencia con el Ejercicio 10 (un solo tipo) es esencial: con un tipo, la
informaci\'on relevante de la pila es solo su \emph{altura} (un n\'umero, $O(\log n)$
bits); con dos o m\'as tipos hay que recordar tambi\'en la \emph{secuencia de tipos},
que es informaci\'on de $\Theta(n)$ bits si no se permite re-leer. La cota
logar\'itmica de la primera parte se logra precisamente \emph{gracias} a poder
re-leer la entrada (acceso aleatorio de solo lectura), recomputando la pila en vez de
almacenarla.
\end{nota}

% ==================================================================
\section{El pal\'indromo requiere $O(n)$ sin retroceder}
% ==================================================================

\begin{teorema}
El problema de reconocer pal\'indromos requiere memoria $\Omega(n)$ (lineal) si la
entrada se lee de izquierda a derecha sin poder volver atr\'as.
\end{teorema}

\begin{idea}
Mismo principio de \emph{distinguibilidad} que en el Ejercicio 11: exhibimos
exponencialmente muchos prefijos que la m\'aquina \emph{tiene} que distinguir, lo que
fuerza un n\'umero exponencial de configuraciones de memoria y por tanto
$\Omega(n)$ bits.
\end{idea}

\begin{proof}
Sea $L_{\text{pal}}=\{w : w=w^{R}\}$ sobre $\{0,1\}$ (donde $w^{R}$ es $w$ invertida).
Supongamos una m\'aquina $M$ que lee la entrada una sola vez, de izquierda a derecha,
con memoria de trabajo de tama\~no $s(n)$.

Fijado $m$, consideramos el conjunto de los $2^{m}$ prefijos $u\in\{0,1\}^{m}$.
Afirmamos que $M$ debe llegar a configuraciones de memoria \emph{distintas} tras leer
dos prefijos distintos $u\neq u'$ (con $|u|=|u'|=m$).

En efecto, supongamos por reducci\'on al absurdo que $M$ alcanza la \emph{misma}
configuraci\'on de memoria tras leer $u$ y tras leer $u'$, con $u\neq u'$. Como
$u\neq u'$, difieren en alguna posici\'on $t$: $u_t\neq u'_t$. Construimos la
continuaci\'on $z=u^{R}$ (el reverso de $u$). Entonces:
\begin{itemize}[nosep]
  \item La cadena $u\,z = u\,u^{R}$ es un pal\'indromo (es de la forma $w w^{R}$),
        as\'i que $M$ debe \textbf{aceptar} $u z$.
  \item La cadena $u'\,z = u'\,u^{R}$ \textbf{no} es un pal\'indromo: su posici\'on
        $t$ (en la primera mitad) vale $u'_t$, mientras que la posici\'on sim\'etrica
        (en la segunda mitad) corresponde a $u_t\neq u'_t$. Luego $M$ debe
        \textbf{rechazar} $u' z$.
\end{itemize}
Pero tras leer $u$ y $u'$ la configuraci\'on de $M$ es la misma \emph{por hip\'otesis},
y como $M$ no puede volver atr\'as, su comportamiento al leer la parte restante $z$ es
id\'entico en ambos casos: aceptar\'ia ambas o rechazar\'ia ambas. Esto contradice que
deba aceptar $uz$ y rechazar $u'z$.

Por tanto las $2^{m}$ configuraciones tras leer los $2^{m}$ prefijos son todas
distintas. El n\'umero de configuraciones de memoria con $s$ celdas (sobre un
alfabeto de trabajo finito de tama\~no $\gamma$, con la cabeza en una de $s$
posiciones y uno de $Q$ estados de control) es a lo sumo
\[
\#\text{config} \le Q\cdot s\cdot \gamma^{s} = 2^{\,O(s)} .
\]
Necesitamos $\#\text{config}\ge 2^{m}$, luego $2^{O(s)}\ge 2^{m}$, de donde
$s = \Omega(m)$. Tomando $m=\lfloor n/2\rfloor$ (cadenas de longitud $n=2m$),
obtenemos
\[
s(n) = \Omega(n).
\]
Es decir, el pal\'indromo requiere memoria \emph{lineal} en el modelo de lectura
\'unica sin retroceso.
\end{proof}

\begin{nota}
Compara con el Ejercicio 5: $\{wcw\}$ \emph{s\'i} est\'a en $\Lclass$ porque con
acceso de \emph{solo lectura con re-posicionamiento} podemos comparar las dos mitades
con punteros ($O(\log n)$). El pal\'indromo $\{ww^{R}\}$ tambi\'en estar\'ia en
$\Lclass$ con re-lectura (comparando $x_i$ con $x_{n+1-i}$ mediante dos punteros). La
cota $\Omega(n)$ de este ejercicio depende \emph{cr\'iticamente} de la restricci\'on
``sin volver atr\'as'': es lo que prohibe usar punteros y obliga a memorizar la
primera mitad entera.
\end{nota}

% ==================================================================
\section{$\NP\neq\SPACE(n)$}
% ==================================================================

\begin{nota}
Este ejercicio es delicado. \textbf{No} se conoce si $\NP\subseteq\SPACE(n)$ o
viceversa: lo que se demuestra es que las dos clases \emph{no pueden coincidir},
exhibiendo una \emph{propiedad estructural} que una tiene y la otra no. Seguimos la
sugerencia del enunciado: $\NP$ es cerrada bajo reducciones en espacio logar\'itmico
(en general, bajo reducciones de muchos-uno polin\'omicas/$\log$-espaciales),
mientras que $\SPACE(n)$ \emph{no} lo es. Una clase cerrada bajo $\log$-reducciones y
otra que no lo es no pueden ser la misma clase. El ingrediente que rompe el cierre de
$\SPACE(n)$ es el \textbf{Teorema de Jerarqu\'ia de Espacio}.
\end{nota}

\begin{definicion}[cierre bajo reducciones en $\Lclass$]
Una clase $\mathcal{D}$ es \emph{cerrada para reducciones en espacio logar\'itmico}
si: siempre que $P_1\reduce P_2$ (reducci\'on muchos-uno calculable en espacio
$\log$) y $P_2\in\mathcal{D}$, se tiene $P_1\in\mathcal{D}$.
\end{definicion}

\subsection*{Paso 1: $\NP$ es cerrada para reducciones en $\Lclass$}

\begin{lema}
Si $P_1\reduce P_2$ (reducci\'on en espacio logar\'itmico) y $P_2\in\NP$, entonces
$P_1\in\NP$.
\end{lema}

\begin{proof}
Sea $f$ la funci\'on de reducci\'on, calculable en espacio $O(\log n)$ (y por tanto en
tiempo polin\'omico, pues una m\'aquina $\log$-espacial tiene a lo sumo
$2^{O(\log n)}=n^{O(1)}$ configuraciones distintas, luego se detiene en tiempo
polin\'omico; en particular $|f(x)|$ es de tama\~no polin\'omico). La reducci\'on
cumple
\[
x\in P_1 \iff f(x)\in P_2 .
\]
Como $P_2\in\NP$, sea $V_2$ su verificador polin\'omico con certificados de tama\~no
$\le q(|y|)$.

Construimos un verificador para $P_1$: dado $x$ y un certificado candidato $c$,
\begin{enumerate}[nosep]
  \item calcular $y=f(x)$ (tiempo polin\'omico);
  \item aceptar syss $V_2(y,c)=1$.
\end{enumerate}
Entonces
\[
x\in P_1 \iff f(x)\in P_2 \iff \exists c\ \big(|c|\le q(|f(x)|)\big):\ V_2(f(x),c)=1,
\]
y como $|f(x)|$ es polin\'omico en $|x|$, el tama\~no del certificado y el tiempo de
verificaci\'on son polin\'omicos en $|x|$. Luego $P_1\in\NP$.
\end{proof}

\subsection*{Paso 2: $\SPACE(n)$ \emph{no} es cerrada para reducciones en $\Lclass$}

La idea es exhibir dos lenguajes $A\reduce B$ con $B\in\SPACE(n)$ pero
$A\notin\SPACE(n)$. La reducci\'on ser\'a un simple ``\emph{padding}'' (relleno) que
agranda la entrada, y la separaci\'on entre ``$A$ dentro'' y ``$A$ fuera'' la
garantiza el teorema de jerarqu\'ia.

\begin{teorema}[Jerarqu\'ia de espacio; Stearns--Hartmanis--Lewis]
Si $s_1(n)=o(s_2(n))$ y $s_2$ es construible en espacio (y $s_2(n)\ge\log n$),
entonces
\[
\SPACE(s_1(n)) \subsetneq \SPACE(s_2(n)),
\]
la inclusi\'on es \emph{estricta}. En particular, para todo $k\ge 1$,
\[
\SPACE(n) \subsetneq \SPACE(n^{2})\subsetneq\cdots\subsetneq\SPACE(n^{k}),
\]
ya que $n = o(n^{2})$, etc.
\end{teorema}

\begin{lema}
$\SPACE(n)$ no es cerrada para reducciones en espacio logar\'itmico.
\end{lema}

\begin{proof}
Por la jerarqu\'ia de espacio, $\SPACE(n)\subsetneq\SPACE(n^{2})$; fijemos un lenguaje
testigo
\[
A\in\SPACE(n^{2})\setminus\SPACE(n) .
\]
($A$ existe por la inclusi\'on estricta: es decidible en espacio cuadr\'atico pero
\emph{no} en espacio lineal.)

Definimos su versi\'on ``rellenada''. A\~nadimos un s\'imbolo nuevo
$\#\notin\Sigma$ y ponemos
\[
B = \{\, x\#^{\,t} \ :\ x\in A,\ t = |x|^{2}-|x| \,\}.
\]
Es decir, a cada $x\in A$ le pegamos suficientes s\'imbolos de relleno para que la
longitud total pase de $|x|$ a $|x|^{2}$. Sea $N=|x\#^{t}|=|x|^{2}$ la longitud de la
cadena rellenada.

\textbf{(i) $B\in\SPACE(n)$.} Con entrada $z$ de longitud $N$: comprobamos que $z$
tiene la forma $x\#^{t}$ con $t=|x|^2-|x|$ (verificable contando, en espacio
$O(\log N)$), extraemos $x$ y ejecutamos el decididor de $A$, que usa espacio
$O(|x|^{2})$. Pero $|x|^{2}=N$, luego ese espacio es $O(N)$, \emph{lineal en la
longitud de la entrada} $z$. Por tanto $B\in\SPACE(N)=\SPACE(n)$. (El relleno
``diluye'' el coste cuadr\'atico en $|x|$ hasta volverlo lineal en $N$.)

\textbf{(ii) $A\reduce B$ en espacio logar\'itmico.} La reducci\'on
$f(x)=x\#^{\,|x|^{2}-|x|}$ se calcula en espacio $O(\log|x|)$: copiar $x$ a la salida
y luego escribir $|x|^{2}-|x|$ s\'imbolos $\#$ llevando solo un contador
($O(\log|x|)$ bits). Claramente
\[
x\in A \iff f(x)=x\#^{t}\in B,
\]
as\'i que es una reducci\'on muchos-uno v\'alida y $\log$-espacial.

\textbf{(iii) Conclusi\'on.} Tenemos $A\reduce B$ con $B\in\SPACE(n)$. Si
$\SPACE(n)$ fuera cerrada para reducciones en $\Lclass$, deber\'iamos concluir
$A\in\SPACE(n)$. Pero elegimos $A\notin\SPACE(n)$. Contradicci\'on. Luego
$\SPACE(n)$ \textbf{no} es cerrada para reducciones en espacio logar\'itmico.
\end{proof}

\subsection*{Conclusi\'on: $\NP\neq\SPACE(n)$}

\begin{proof}
Combinamos los dos pasos:
\begin{itemize}[nosep]
  \item $\NP$ \emph{s\'i} es cerrada para reducciones en espacio logar\'itmico
        (Paso 1).
  \item $\SPACE(n)$ \emph{no} es cerrada para reducciones en espacio logar\'itmico
        (Paso 2).
\end{itemize}
El cierre para $\log$-reducciones es una propiedad que una clase tiene o no tiene; no
puede ser que dos clases \emph{iguales} difieran en ella. Como $\NP$ la posee y
$\SPACE(n)$ no, forzosamente
\[
\boxed{\ \NP\neq\SPACE(n).\ }
\]
\end{proof}

\begin{nota}[por qu\'e este argumento no decide la inclusi\'on]
El razonamiento prueba la \emph{desigualdad} de clases sin decir cu\'al contiene a
cu\'al --- de hecho hoy se desconoce si $\NP\subseteq\SPACE(n)$. La t\'ecnica
(``una clase es cerrada bajo cierto tipo de reducci\'on y la otra no, luego son
distintas'') es un patr\'on cl\'asico para separar clases. El \'unico ingrediente no
elemental es el Teorema de Jerarqu\'ia de Espacio, que proporciona el lenguaje
$A\in\SPACE(n^2)\setminus\SPACE(n)$; este teorema se demuestra por
\emph{diagonalizaci\'on} (se construye una m\'aquina que en espacio $n^2$ contradice a
toda m\'aquina de espacio lineal). La maniobra de \emph{padding} ``traslada'' esa
separaci\'on a un par de lenguajes relacionados por una reducci\'on trivial,
exhibiendo el fallo de cierre.
\end{nota}

\vspace{6mm}
\hrule
\vspace{3mm}
\begin{center}\small
Fin de la Relaci\'on 4.\\[1mm]
Documento elaborado con fines de estudio. Los resultados son est\'andar en teor\'ia
de la complejidad (Sipser; Papadimitriou; Garey--Johnson); las demostraciones siguen
los esquemas habituales de esos textos y la terminolog\'ia de la asignatura.
\end{center}

\end{document}
